% =============================================================================
\chapter{Hipóteses e Arquitetura Proposta}
% =============================================================================

\section{Hipótese principal}

A hipótese principal é que uma política MORL condicionada por preferências
pode selecionar métodos de decomposição HTN de forma mais adaptativa que
prioridades estáticas ou heurísticas de pesos fixos, mantendo a validade
semântica garantida pelo HTN.

\section{Hipóteses secundárias}

\begin{description}
\item[H1 --- desempenho adaptativo:] a arquitetura completa obterá melhor
desempenho multiobjetivo sob mudanças ambientais que um HTN com ordem fixa de
métodos.

\item[H2 --- generalização de preferências:] uma política condicionada por
vetores de preferência alterará suas escolhas de maneira consistente quando a
importância relativa dos objetivos mudar, inclusive para combinações não
observadas diretamente no treinamento.

\item[H3 --- validade semântica:] como o HTN filtra os métodos antes da decisão
MORL, a arquitetura completa não selecionará métodos que violem pré-condições
ou restrições definidas pelo domínio.

\item[H4 --- interpretabilidade:] cada escolha poderá ser explicada pela tarefa
composta atual, pelos métodos aplicáveis e rejeitados, pelo vetor de
preferências, pelos valores multiobjetivo estimados e pela utilidade
escalarizada final.

\item[H5 --- redução do espaço de decisão:] restringir o MORL aos métodos HTN
aplicáveis produzirá um espaço de ação menor e semanticamente mais significativo
que o aprendizado direto de ações primitivas.
\end{description}

\section{Componentes da arquitetura}

\subsection{HTN: validade e estrutura}

O HTN define tarefas compostas e primitivas, métodos, pré-condições, requisitos
de recursos e restrições narrativas ou de gameplay. Para cada decisão, ele
produz $M_{\mathrm{valid}}(s_t,\tau_t)$. Métodos impossíveis ou proibidos são
eliminados antes da consulta ao modelo aprendido.

\subsection{Gerador de Pesos Orientado a Objetivos}

O GOWG produz $\mathbf{w}_t$ a partir do objetivo, estado, personalidade e
contexto. Exemplos de componentes são sobrevivência, missão, proteção de
aliados, conservação de recursos, furtividade e velocidade. Regras como
``aumentar o peso de sobrevivência quando a saúde diminui'' serão explicitadas
e versionadas para garantir reprodutibilidade.

\subsection{MORL: seleção de método}

O MORL recebe estado, tarefa composta, vetor de preferência e máscara de
métodos válidos. Sua ação é escolher um método HTN. Portanto, a formulação
corrige a ambiguidade da hipótese inicial: o MORL não é a camada motora e não
substitui os controladores de execução. Ele decide \emph{como decompor} a tarefa
corrente entre alternativas semanticamente permitidas.

\section{Exemplo de decisão}

Para a tarefa composta \texttt{HandleThreat}, o domínio pode definir os métodos
\texttt{Direct\-Engagement}, \texttt{Seek\-Cover},
\texttt{Protect\-Ally}, \texttt{Retreat} e \texttt{Flank}. Se não
houver rota lateral, \texttt{Flank} é removido. Se não houver aliado,
\texttt{Protect\-Ally} é removido. O MORL compara apenas os métodos restantes sob
$\mathbf{w}_t$.

Uma explicação de decisão poderá assumir a forma: o agente selecionou
\texttt{Protect\-Ally} porque o método era aplicável e apresentou a maior
utilidade estimada sob um vetor no qual proteção do aliado e sobrevivência
possuíam os maiores pesos.

\section{Ciclo online}

\begin{enumerate}
\item o HTN identifica a tarefa composta corrente;
\item o HTN calcula os métodos aplicáveis;
\item o GOWG gera o vetor de preferências;
\item o MORL seleciona um método válido;
\item o método é decomposto em subtarefas;
\item controladores tradicionais executam as tarefas primitivas;
\item o ambiente fornece novo estado e recompensa vetorial;
\item o ciclo reinicia na próxima decisão hierárquica.
\end{enumerate}
